\documentclass{ximera}[font=12pt]


\title{Adding Integers: Different Signs}   % Use xmlatex all -s to compile when SERVE refuses to work.
\author{Kelly Stady}
\license{CC: 0}         % replace with an appropriate license, or set it in xmPreamble

\begin{document}

\begin{abstract}

\end{abstract}

\maketitle
\label{xim:AddIntDiffSActivity}


\section*{Adding Integers with Different Signs}

When adding integers with different signs, we can think of the process as a
\textbf{battle between positive and negative numbers}, where the winner is the sign of
the answer.
/

\textbf{Battle 1:}  \ If you have \$30 and spend \$20 to see a movie, your net balance is \$10. 

\[ 30 + (-20) = 10 \]

\hspace{1in}  \textbf{The answer is positive, so the positive team win!}


\textbf{Battle 2:} \ If you have \$30 and a credit card debt of \$100, your net balance is a debt of \$70.


\[ 30 + (-100) = -70 \]

\hspace{1in}  \textbf{The answer is negative, so the negative team win!}

Why did the positive team win the first battle and lose the second battle?

\textbf{Answer:}  The positive team won the first battle because the asset was greater than the debt, but 
lost the second because the debt was greater than the asset. 


In any number battle, the winner is always the team with the \textbf{greatest strength}, regardless of its sign.  
Mathematically, we find the \textbf{absolute value} of a number to determine its \textbf{strength} or \textbf{size}.  


\textbf{Battle 1:} \ $30 + (-20) = 10$  
\begin{itemize}
    \item[] The strength of the \$30 asset is $|30| = 30$.
    \item[] The strength of the \$20 debt is $|-20| = 20$.
\end{itemize}


Since $|30|$ is larger than $|-20|,$ the \textbf{positive} team was stronger and \textbf{won} the battle.


\textbf{Battle 2:} \ $30 + (-100) = -70$
\begin{itemize}
    \item[] The strength of the \$30 asset is $|30| = 30$.
    \item[] The strength of the \$100 debt is $|-100| = 100$.
\end{itemize}

Since $|-100|$ is larger than $|30|,$ the \textbf{negative} team was stronger and \textbf{won} the battle.

We are now ready for a formal mathematical rule that can be applied to any problem, not just those involving money.  
As we noticed in our earlier examples, adding numbers with different signs is a bit of a paradox: to find the sum, 
we actually have to subtract.  From the number battles, we know that the sign of the answer depends on which number 
has the greatest strength. To turn these observations into a formal rule, we will use absolute value.


\begin{procedure}[Steps for Adding Integers with Different Signs]

\begin{enumerate} 
    \item Find the \textbf{absolute value} or \textbf{strength} of each number.

     \item Subtract the smaller absolute value from the larger absolute value.

     \item Write the answer with the sign of the number with the \textbf{larger absolute value}.
\end{enumerate}
\end{procedure}


\begin{example} Let's use our three-step process to find the sum $-50 + 30.$

\begin{description}[style=multiline, leftmargin=5em, font=\bfseries]
    \item[Step 1] \textbf{Find the Absolute Values (The Strengths):} \\
    Find the strength of each number.
    \[ |-50| = 50 \quad \text{and} \quad |30| = 30 \]

    \item[Step 2] \textbf{Subtract the Absolute Values (The Battle):} \\
    Since we are adding numbers with different signs, the numbers ``battle" by finding the difference between their strengths.
    \[ 50 - 30 = 20 \]

    \item[Step 3] \textbf{Declare the Winner (The Sign):} \\
    Compare the absolute values from the first step. Since $|-50|$ is greater than $|30|,$ the \textbf{negative team}
     is stronger and \textbf{wins} the battle!  
    \[ -50 + 30 = -20 \]
\end{description}

\end{example}


\bigskip

\begin{problem}
    Apply the steps to find the sum $15 + (-40).$

    ~ \\[3em]

    \textbf{Step 1:} \textbf{Find the Absolute Values (The Strengths)}:
    \[ |15| = \answer{15} \quad \text{and} \quad |-40| = \answer{40} \]

    
    \textbf{Step 2:} \textbf{Subtract the Absolute Values (The Battle)}:
    \[ 40 - 15 = \answer{25} \]


    \textbf{Step 3:} \textbf{Declare the Winner (The Sign):}

    Compare the absolute values: is $|15|$ or $|-40|$ larger?
    
    \begin{multipleChoice}
        \choice{The positive team wins, since $|15|$ is larger.}
        \choice[correct]{The negative team wins, since $|-40|$ is larger.}
    \end{multipleChoice}

    ~ \\[2em]

    \textbf{Final Answer:} $15 + (-40) = \answer{-25}$

\end{problem}


\begin{problem}
Find the following sums.

\[ -12 + 4 = \begin{prompt} \answer{}  \end{prompt} \]

\[ 12 + (-4) = \begin{prompt} \answer{}  \end{prompt} \]

\[ 32 + -20  = \begin{prompt} \answer{}  \end{prompt} \]

\[ 100 + (-600) = \begin{prompt} \answer{}  \end{prompt} \]


\end{problem}

~ \\[2em]

\begin{problem}
Atoms are the tiny building blocks of all matter.  An atom consists of protons, electrons and neutrons.    
Each proton has a positive charge ($+1$), each electron has a negative charge ($-1$) and neutrons have no charge (0).  
Chemists find the charge of an atom by adding the charges of its electrons and protons.  Write the addition 
problem to find the charge of an atom that has eight electrons and three protons.  Then find the charge of the atom.

~ \\[2em] 

\begin{center}
\begin{tikzpicture}[scale=0.8]
    % Protons: 3 arranged in a triangle
    \foreach \a in {90, 210, 330} {
        \fill[orange!85!black] (\a:0.35) circle (0.3);
        % Adjusted Plus Sign (slightly thinner than before)
        \fill[white] (\a:0.35) +(-0.16,-0.03) rectangle +(0.16,0.03);
        \fill[white] (\a:0.35) +(-0.03,-0.16) rectangle +(0.03,0.16);
    }

    % Inner shell: 2 electrons
    \draw[gray!50, thick] (0,0) circle (1.3);
    \foreach \a in {45, 225} {
        \fill[glaucous] (\a:1.3) circle (0.25);
        % Adjusted Minus Sign (0.03 thickness instead of 0.05)
        \fill[white] (\a:1.3) +(-0.14,-0.03) rectangle +(0.14,0.03);
    }

    % Outer shell: 6 electrons
    \draw[gray!60, thick] (0,0) circle (2.4);
    \foreach \a in {0, 60, 120, 180, 240, 300} {
        \fill[glaucous] (\a:2.4) circle (0.25);
        % Adjusted Minus Sign
        \fill[white] (\a:2.4) +(-0.14,-0.03) rectangle +(0.14,0.03);
    }
\end{tikzpicture}

\xmalt{A model of the atom. The nucleus has 3 large orange protons with white plus signs. Two gray orbits surround it: 
an inner orbit with 2 blue electrons and a slightly larger outer orbit with 6 blue electrons. All electrons have white minus signs.}


\end{center}

~ \\[2em] 

\begin{multipleChoice}
\choice[correct]{$-8 + 3$}
\choice{$-8 + (-3)$}
\choice{$8 + (-3)$}
\choice{$8 + 3$}
\end{multipleChoice}

Charge $= \answer{-5}$
\end{problem}




\end{document}